%%%%%%%%%%%%%%%%%%%%%%%%%%%%%%%%%%%%%%%
% General two-page resume
% Deedy two-column layout, split manually across two pages
% (minipages cannot break across pages).
% Build from the repo root:  .\build_variants.ps1
%%%%%%%%%%%%%%%%%%%%%%%%%%%%%%%%%%%%%%%

\documentclass[]{deedy-resume-openfont}
\usepackage{fancyhdr}
\geometry{a4paper,left=1.0cm,right=1.0cm,top=0.6cm,bottom=1.0cm}

\pagestyle{fancy}
\fancyhf{}
\renewcommand{\headrulewidth}{0pt}
\pagefooter{Mikołaj Mieszko Charchuta \textbullet{} Resume}
\fixlocationgap

\begin{document}

\renewcommand{\sectionsep}{\vspace{7pt}}

\resumephoto[19mm]{6mm}{3.5mm}

%%%%%%%%%%%%%%%%%%%%%%%%%%%%%%%%%%%%%%
%     TITLE NAME
%%%%%%%%%%%%%%%%%%%%%%%%%%%%%%%%%%%%%%
\namesection{Mikołaj Mieszko}{Charchuta}{ \urlstyle{same}\href{https://www.github.com/mmcharchuta}{github.com/mmcharchuta} | \href{https://www.linkedin.com/in/mmcharchuta/}{linkedin.com/in/mmcharchuta/}\\ +48 730 359 228 | \href{mailto:mmcharchuta@gmail.com}{mmcharchuta@gmail.com} | Vienna, Austria (open to remote, hybrid, relocation)
}

%%%%%%%%%%%%%%%%%%%%%%%%%%%%%%%%%%%%%%
%     PROFILE
%%%%%%%%%%%%%%%%%%%%%%%%%%%%%%%%%%%%%%

\vspace{3pt}
\noindent{\small \textbf{Bioinformatician turning exploratory genomics into reproducible, production-grade pipelines} (Nextflow, Snakemake, Docker, Conda), specialising in chromosome conformation capture (Meta Hi-C, 3C-seq) and machine learning on bacterial pangenomes.
Bridging academia and industry: Meta Hi-C for phage-therapy R\&D at Phagos, an independently grant-funded project on the 3D genome of \textit{B. thuringiensis}, and prior R\&D in personalised medicine and cellular agriculture.
Three publications and invited teaching on reproducibility; EU remote-first, open to B2B contracts and industrial PhD (CIFRE) partnerships.}
\vspace{3pt}

%%%%%%%%%%%%%%%%%%%%%%%%%%%%%%%%%%%%%%
%     PAGE 1 -- COLUMN ONE
%%%%%%%%%%%%%%%%%%%%%%%%%%%%%%%%%%%%%%

\begin{minipage}[t]{0.33\textwidth}

\section{Skills}
\subsection{Reproducibility}
Nextflow \textbullet{} Snakemake \textbullet{} Docker \\
Conda \textbullet{} Git \textbullet{} CI/CD \\
Slurm \textbullet{} AWS
\sectionsep

\subsection{Programming \& ML}
R \textbullet{} Python \textbullet{} Bash \textbullet{} C++ \\
scikit-learn \textbullet{} tidymodels \\
R package development (NaïveR)
\sectionsep

\subsection{Genomics}
3D genomics (Hi-C, Meta Hi-C, 3C-seq) \\
Pangenomics \textbullet{} WGS \textbullet{} WGBS \\
samtools \textbullet{} GATK \textbullet{} IGV \textbullet{} ChIP-seq
\sectionsep

\subsection{Phages \& Proteins}
Phage--host association (Meta Hi-C) \\
Toxin \& gene mining in pangenomes \\
Protein structure (PyMOL) \\
Proteomics (PEAKS Studio, ProteinPilot)
\sectionsep

\subsection{Transcriptomics}
DESeq2 \textbullet{} rMATS \textbullet{} rnaSPAdes
\sectionsep

\subsection{Metagenomics}
QIIME2 \textbullet{} mothur \textbullet{} NaïveR
\sectionsep

\subsection{Sequencing \& Wet Lab}
Nanopore \textbullet{} Illumina \textbullet{} Ion Torrent \\
Hi-C library prep \textbullet{} qPCR \\
DNA isolation \textbullet{} animal work
\sectionsep

\section{Education}

\subsection{BOKU, Vienna}
\descript{MSc Bioinformatics (Erasmus)}
\location{Sep 2026 -- Jul 2027}
{\footnotesize Academic exchange; thesis project with Prof. Christian Stauffer.}\\
\sectionsep

\subsection{AMU, Poznan}
\descript{MSc in Bioinformatics}
\location{Oct 2025 -- Jul 2027}
{\footnotesize Thesis: \textit{In silico analysis of Cry9 toxins and their genetic determinants in the Bacillus thuringiensis pangenome} (first-author article in preparation).}\\
\sectionsep

\subsection{AMU, Poznan}
\descript{BSc Bioinformatics \& Biotechnology}
\location{Oct 2021 -- Sep 2025}
{\footnotesize Thesis: \textit{naïve methods for taxonomic classification from marker genes}, released as the \href{https://github.com/mmcharchuta/NaiveR}{NaïveR R package}.}\\
\sectionsep

\end{minipage}
\hfill
\begin{minipage}[t]{0.66\textwidth}

%%%%%%%%%%%%%%%%%%%%%%%%%%%%%%%%%%%%%%
%     PAGE 1 -- EXPERIENCE
%%%%%%%%%%%%%%%%%%%%%%%%%%%%%%%%%%%%%%

\section{Experience}

\runsubsection{Phagos}
\descript{| Bioinformatician}
\location{Jul 2026 -- Present | Remote / Paris, FR}
\begin{tightemize}
\item Deliver phage--host association calls into the therapeutic discovery pipeline by analysing Meta Hi-C data end-to-end, from raw reads to ranked pairings.
\item Made analyses re-runnable and free of environment drift by rebuilding ad-hoc scripts as containerised Nextflow workflows (Docker, Conda) used across the R\&D team.
\item Guide two research projects and joint analyses with external academic partners (incl. Institut Pasteur), aligning experimental design with computational feasibility.
\end{tightemize}
\sectionsep

\runsubsection{AMU, Dept. of Microbiology}
\descript{| Contract Researcher}
\location{Sep 2025 -- Jan 2028 | Remote / Poznań, PL}
\begin{tightemize}
\item Won competitive Study@Research funding (app. nr. 201/34/UAM/0074) to run an independent project on the 3D genome of \textit{B. thuringiensis}, owning Hi-C library prep, sequencing and analysis solo.
\item Produced the dataset behind a first-author manuscript by mining the \textit{B. thuringiensis} pangenome for Cry9 toxin determinants.
\item Work on genomes of insect pests and biopesticidal microbes, combining comparative genomics with predictive protein design.
\item Standardised the group's Hi-C/3C-seq analysis with reusable containerised workflows for contact-map QC and scaffolding.
\end{tightemize}
\sectionsep

\runsubsection{Gesellschaft für Bioanalytik}
\descript{| Scientist}
\location{Jul 2024 -- Jun 2026 | Poznań, PL}
\begin{tightemize}
\item Delivered allergen and biosafety analyses to commercial turnaround and QA standards.
\item Covered the whole chain from field sampling through lab work to computational analysis and client reporting.
\end{tightemize}
\sectionsep

\runsubsection{Bene Meat Technologies}
\descript{| Bioinformatician}
\location{Mar 2024 -- Jul 2024 | Prague, CZ}
\begin{tightemize}
\item Provided end-to-end bioinformatics support to cultivated-meat cell-line R\&D.
\item Ran transcriptomics (DESeq2, rMATS) and nanopore analyses, and automated bioreactor workflows.
\end{tightemize}
\sectionsep

\runsubsection{Leiden University Medical Center}
\descript{| Bioinformatician}
\location{Jun 2023 -- Oct 2023 | Leiden, NL}
\begin{tightemize}
\item Contributed to two peer-reviewed publications by building the pipelines and figures for a murine MASLD study in the Endocrinology Department.
\item Provided bioinformatics support to the lab and conducted animal research (mice, tissue and organ work).
\end{tightemize}
\sectionsep

\runsubsection{MNM Diagnostics}
\descript{| Genome Analyst}
\location{Feb 2022 -- Jul 2022 | Poznań, PL}
\begin{tightemize}
\item Interpreted clinical whole-genome data for personalised oncology, reporting variants to clinicians.
\item Supported the machine-learning team with feature extraction from methylation (WGBS) data, the data lake and explainable-AI (XAI) work.
\item Automated tasks and data acquisition, and worked with biobanks under strict patient-data safety rules.
\end{tightemize}

\end{minipage}

%%%%%%%%%%%%%%%%%%%%%%%%%%%%%%%%%%%%%%
%     PAGE 2
%%%%%%%%%%%%%%%%%%%%%%%%%%%%%%%%%%%%%%
\newpage

\begin{minipage}[t]{0.33\textwidth}

\section{Research Exchanges}

\subsection{Charles Univ., Prague}
\descript{Erasmus+}
\location{Winter 2022/23 \textbullet{} Summer 2023/24}
{\footnotesize \textbullet{} \textbf{Institute of Microbiology, CAS}: stable isotope probing of polar soil microbes for carbon tracking; meta-omics. \\
\textbullet{} \textbf{BIOCEV}: alternative splicing analysis (rMATS) on human RNA-seq data. \\
\textbullet{} CryoEM of cultivated-meat proteins.}\\
\sectionsep

\subsection{Univ. of Gdansk \& GUMed}
\descript{National exchange}
\location{Winter 2023/24}
{\footnotesize \textbullet{} Metabolomics and proteomics of bioprotective synthetic communities (SynComs). \\
\textbullet{} Phage genomics.}\\
\sectionsep

\section{Grants \& Awards}
{\footnotesize
\textbullet{} Study@Research, AMU (2026--27): 3D genome of \textit{B. thuringiensis} \\
\textbullet{} SPUN international grant (2022--24): ML for soil biodiversity hotspots \\
\textbullet{} AMU grants for talented students: DynaMAP (2023--25), Trees of the Future (2021--23) \\
\textbullet{} Best poster, national metagenomics symposium}
\sectionsep

\section{Languages}
English (C2) \textbullet{} Polish (Native)  \\
German (B2) \textbullet{} French (B2) \\
Czech (A2) \textbullet{} Hungarian (A2)
\sectionsep

\section{Interests}
{\footnotesize Beekeeping \& fungiculture \textbullet{} bikepacking \textbullet{} vegan cooking}

\end{minipage}
\hfill
\begin{minipage}[t]{0.66\textwidth}

\section{Research Projects}

\runsubsection{DynaMAP nanopore profiling}
\descript{| Team lead, funded BSc thesis project}
\location{Oct 2023 -- Oct 2025 | AMU, Poznań}
\begin{tightemize}
\item Led a student team benchmarking a new nanopore-based microbiome profiling technique against established standards, showing its potential to lower sequencing costs.
\end{tightemize}
\sectionsep

\runsubsection{SPUN global soil mapping}
\descript{| International grant}
\location{Nov 2022 -- Apr 2024 | Illumina}
\begin{tightemize}
\item Predicted soil biodiversity hotspots with machine learning for a global mapping project (\href{https://youtu.be/6frg_7ykXvc?t=128}{field video}).
\end{tightemize}
\sectionsep

\runsubsection{Trees of the Future}
\descript{| Team grant}
\location{Oct 2021 -- Oct 2023 | Ion Torrent}
\begin{tightemize}
\item Identified interaction hubs between soil microorganisms; awarded best poster at the national metagenomics symposium; article in preparation.
\end{tightemize}
\sectionsep

\runsubsection{NaïveR}
\descript{| Open-source R package}
\location{\href{https://github.com/mmcharchuta/NaiveR}{github.com/mmcharchuta/NaiveR}}
\begin{tightemize}
\item Naïve methods for taxonomic classification from marker genes, released from my BSc thesis.
\end{tightemize}
\sectionsep

\section{Teaching, Talks \& Community}

\runsubsection{Danube Int. School Vienna}
\descript{| Lab Assistant}
\location{Jun 2026 -- Aug 2028 | Vienna, AT}
\begin{tightemize}
\item Design and run hands-on science labs; mentor students through independent research projects.
\end{tightemize}
\sectionsep

\runsubsection{CRABS 2026 / Health Intel.}
\descript{| Invited Speaker \& Workshop Lead}
\location{Aug 2026 | Kyiv, UA}
\begin{tightemize}
\item Lecture on \textit{Reproducibility in Computational Genomics}; two Nextflow pipeline workshops; talk on Meta Hi-C \& ML for bacterial genomes.
\end{tightemize}
\sectionsep

\runsubsection{Bioinformatics Community PL}
\descript{| Co-organiser}
\location{2026 -- Present}
\begin{tightemize}
\item Event coordination, website and open-science advocacy.
\end{tightemize}
\sectionsep

\section{Publications}
{\footnotesize
\textbf{In silico analysis of Cry9 toxins and their genetic determinants in the \textit{Bacillus thuringiensis} pangenome.} \textbf{Charchuta, M. M.} (first author), Baranek, J., et al. \textit{In preparation, 2026.}\\[2pt]
\textbf{Combination of dietary fiber and exercise training improves fat loss in mice but does not ameliorate MASLD more than exercise alone.} Kovynev, A.; \textbf{Charchuta, M. M.}; et al. \textit{Am. J. Physiol. Gastrointest. Liver Physiol.}, 2025.\\[2pt]
\textbf{WED-517 Oatmeal a day does not keep the doctor away: inulin and exercise impact on MASLD amelioration and gut-liver crosstalk.} Kovynev, A.; \textbf{Charchuta, M.}; et al. \textit{J. Hepatol.}, 2024.
}

\end{minipage}

\vspace*{\fill}
\noindent{\tiny I consent to the use of my personal data for recruitment purposes, in line with the GDPR (EU) 2016/679. My information will be stored and processed solely for the duration of the hiring process.}

\end{document}
